% !TEX root = ../CerdaMS.tex

\section{Taylor approximation for Ohana trees}\label{sec:TE4MTs}

In its original formulation, Ehrhard and Regnier's Taylor expansion translates a \lam-term as a power series of iterated derivatives~\cite{Ehrhard.Reg.06}.
We now introduce a \emph{qualitative} Taylor expansion~\cite{ManzonettoP11} specifically designed for the \lamI-calculus, having as target the \lamI-resource calculus.

\subsection{The Taylor approximation}

Intuitively, a qualitative Taylor expansion should associate each term $M \in \LamI(X)$ with a \emph{set} of resource approximants $\TE M \subseteq \rterms(X)$.
Therefore the codomain of $\TE -$ should be $\bigdunion_{X \subf \Var} \parts{\rterms(X)}$, and what we actually define is $\TE M \eqdef (X, \TE* M)$ with $\TE* M \subseteq \rterms(X)$.
Notice that, whereas all previous constructions of disjoint unions
$\bigdunion_{X \subf \Var}$ were formed from genuinely disjoint sets,
this is no longer the case here, as $\emptyset \in \parts{\rterms(X)}$, for all $X$.
\begin{notation}\mbox{}
\bsub
\item For $X \subf \Var$ and $\cX \in \parts{\rterms(X)}$,
	we write $\fv(X,\cX) \eqdef X$.
\item We let $\mathord{\TElt}$ denote the order relation
	on $\bigdunion_{X \subf \Var} \parts{\rterms(X)}$
	such that $(X,\cX) \TElt (Y,\cY)$
	whenever $X = Y$ and $\cX \subseteq \cY$.
	We write $\TElub$ for the corresponding least upper bounds.
\item We write $\rt s \TEin (X,\cX)$ to mean that $\rt s \in \cX$.
\item Given $\cX \in \parts{\rterms(X)}$, 
	$\cY \in \parts{\rterms*(Y)} \setminus \set\emptyset $
	and $x \in X$, we write:
%	\begin{align*}
%	\lam x.\cX & \eqdef \set{\lam x.\rt s}[\rt s\in\cX] &
%	\cX\cY & \eqdef 
%		\set{\rt s\rb t}[\rt s \in \cX,\ \rb t \in \cY] \\
%	\lam x.(X,\cX) & \eqdef (X \setminus \set x, \lam x.\cX) &
%	(X,\cX) \cY & \eqdef (X \cup Y, \cX\cY).
%	\end{align*}
	\[	\lam x.(X,\cX) \eqdef (X \setminus \set x, \set{\lam x.\rt s}[\rt s\in\cX]), \qquad
		(X,\cX) \cY \eqdef (X \cup Y, \set{\rt s\rb t}[\rt s \in \cX,\ \rb t \in \cY]). 
	\]
\esub
\end{notation}

\begin{definition}\mbox{}
\bsub
\item The \defemph{Taylor expansion} $\TE{M} \in \bigdunion_{X \subf \Var} \parts{\rterms(X)}$ of a \lamI-term $M$, is defined together with $\TE** M \subset \rterms*(\fv(M))$ by mutual induction:
	\begin{gather*}
		\begin{aligned}
		\TE{x} & \eqdef (\set{x}, \set{x}), &
		\TE{\lam x.M} & \eqdef \lam x.\TE M, &
		\TE{MN} &\eqdef  \TE M \TE** N, 
%		\\
%		\TE*{x} & \eqdef \set{x} &
%		\TE*{\lam x.M} & \eqdef \lam x.\TE* M &
%		\TE*{MN} \eqdef & \TE* M \TE** N
		\end{aligned} \\
		\TE**{M} \eqdef \set{1_{\fv(M)}} \cup 
			\set{ [t_0,\dots,t_n] } [ n\in\Nat,\ t_0,\dots,t_n\TEin\TE M ]. 
	\end{gather*}
\item The above definition is extended to Ohana trees by setting,
	for all $M \in \LamI$,
	\[	\TMT M \eqdef \bigTElub_{A \in \Appof{M}} \TE{A},
%		\qquad
%		\TMT* M \eqdef \bigcup_{A \in \Appof{M}} \TE*{A}
		\qquad \text{together with }
		\left\{ \begin{array}{r!{\eqdef}l}
			\TE{\bot_X} & (X,\emptyset), \\ 
			\TE**{\bot_X} & \set{1_X}.\\
		\end{array} \right.
	\]
%	together with $\TE{\bot_X} \eqdef (X,\emptyset)$,
%	$\TE*{\bot_X} \eqdef \emptyset$
%	and $\TE**{\bot_X} \eqdef \set{1_X}$.
\esub
\end{definition}

\begin{remark} \label{rem:fv-of-taylor}
	For all $M$, we have $\fv(\TE M) = \fv(M)$.
%	Notice that for all $M$, $\TE M = (\fv(M), \TE* M)$ 
%	for some $\TE* M \in \parts{\rterms(\fv(M))}$.
	Similarly, $\fv(\TMT M) = \fv(M)$,
%	Similarly, $\TMT M  = (\fv(M), \TMT* M)$ 
%	for some $\TMT* M \in \parts{\rterms(\fv(M))}$,
	due to \cref{rem:appof-M-fv-M} and the way
	we ordered $\bigdunion_{X \subf \Var} \parts{\rterms(X)}$.
\end{remark}

\begin{example}\mbox{}
	\bsub
	\item $\TE{\comb I} = (\emptyset, \set{\comb I})$,
		$\TE{\comb D} = (\emptyset, \set{\comb D_n}[n \in \Nat])$.
	\item $\TMT{\comb Yf} = (\set f, \cX_{\comb Yf})$ and
		$\TMT{\Bible[l]f} = (\set{f,l}, \cX_{\Bible[l]f})$,
		where the sets of approximants can be described 
		as the smallest (in fact, unique) subsets of $\rterms$ such that:
		\begin{align*}
			\cX_{\comb Yf} & = \set{ f1_{\set f} } \cup 
				\set{ f[\rt t_0,\dots,\rt t_n] }
				[n\in\Nat,\ \rt t_0,\dots,\rt t_n \in \cX_{\comb Yf} ], \\
			\cX_{\Bible[l]f} & = \set{ f1_{\set{f,l}} } \cup 
				\set{ f[\rt t_0,\dots,\rt t_n] }
				[n\in\Nat,\ \rt t_0,\dots,\rt t_n \in \cX_{\Bible[l]f} ].
		\end{align*}
	\esub
\end{example}

We now describe how resource reduction acts on Taylor expansions, \emph{i.e.\ }on potentially infinite sets of resource expressions.

\begin{notation} \label{not:support}\mbox{}
\bsub
\item Given $\rs s \in \frsums[X]**$,
	we denote by $\supp s \in \bigdunion_{X \subf \Var} \parts{\rterms(X)}$ 
	its \defemph{support}, defined by
	$\supp*{\esum X} \eqdef (X,\emptyset)$ and
	$\supp*{\rt s + \rs t} \eqdef (X,\set{\rt s}) \TElub \supp t$.
	Notice that $\rt s \TEin \supp t$ whenever $\rt s$ 
	bears a non-zero coefficient in $\rs t$ (i.e.\ $\rs t = s + \rs t'$, for some $\rs t'$).
\item Given $\cX \in \parts{\rterms**(X)}$, we write 
	$\nf(X, \cX) \eqdef \bigTElub_{\rt s \in \cX} \supp*{ \nf(\rt s) }$.
	Notice that $\fv(\nf(X, \cX)) = X$, and that
	$\nf(X, \cX) = (X, \emptyset)$ if and only if
	$\rt s \rreds \esum X$, for all $\rt s \in \cX$.
\esub
\end{notation}

This allows us to state the following theorem,
adapting \cite[Theorem~2]{Ehrhard.Reg.06}.

\begin{theorem}[Commutation] \label{the:commutation}
	For all $M \in \LamI$, $\NFT{M} = \TMT{M}$.
\end{theorem}

We outline a proof in the following sequence of lemmas,
similar to the one in \cite{Barbarossa.Man.19};
alternatively, the variants of \cite{Ehrhard.Reg.06,Ehrhard.Reg.08},
\cite{OlimpieriVaux22} and \cite{Cerda.24,CerdaVaux23}
could also be adapted.

\newcounter{lem:rred-conserv-bred}
\setcounter{lem:rred-conserv-bred}{\thetheorem}
\begin{lemma}
\label{lem:rred-conserv-bred}
	For all resource sums $\rs s$ such that $\supp s \TElt \TE{M}$,
	there exists a reduction $M \breds N$
	such that $\supp*{\nf(\rs s)} \TElt \TE{N}$.
\end{lemma}

	\begin{proof}
	By induction on the length of the longest reduction 
	$\rs s \rreds \nf(\rs s)$.
	If it is $0$, take $N \eqdef M$.
	Otherwise, the longest reduction is
	$\rs s = \rt t + \rs u \rred \rs t' + \rs u \rreds \nf(\rs s)$.
	By firing all redexes in $M$ and $\rs u$ at the same position
	as the redex fired in $\rt t \rred \rs t'$
	(formally we do this by induction on this reduction),
	we obtain $M \bred M'$ and $\rs u \rreds \rs u'$ such that
	$\supp*{\rs t' + \rs u'} \TElt \TE{M'}$.
	By confluence, $\rs t' + \rs u' \rreds \nf(\rs s)$.
	We conclude by induction on this (shorter) reduction.
	\end{proof}

\newcounter{lem:normal-approximants}
\setcounter{lem:normal-approximants}{\thetheorem}
\begin{lemma}
\label{lem:normal-approximants}
	If $\rt t \TEin \TE N$ is in $\resource$-nf, then
	there exists $A \in \Appof{N}$ such that $\rt t \TEin \TE A$.
\end{lemma}

	\begin{proof}
	By induction on $t$. An $\resource$-nf is also a hnf, hence
	$\rt t = \lam \vec{x}.y\rb t_1\cdots\rb t_k$ and
	$N = \lam \vec{x}.yN_1\cdots N_k$.
	For $1 \leq i \leq k$, $\rb t_i \in \TE**{N_i}$.
	If $\rb t_i = 1_{\fv(N_i)}$ define $A_i \eqdef \bot_{\fv(N_i)}$.
	If $\rb t_i = [t_i^1,\dots,t_i^n]$,
	each $t_i^j$ is in $\resource$-nf, by induction
	there is $A_i^j \in \Appof{N_i}$, $t_i^j \TEin \TE{A_i^j}$.
	Define $A_i \eqdef \bigMTlub_{j=1}^n A_i^j \in \Appof{N_i}$.
	Finally, $A \eqdef \lam \vec{x}.yA_1\cdots A_k \in \Appof{N}$
	and $t \TEin \TE A$.
	\end{proof}

\newcounter{lem:TE-monotone}
\setcounter{lem:TE-monotone}{\thetheorem}
\begin{lemma}[Monotonicity of $\TE{-}$]\mbox{}
\label{lem:TE-monotone}
	\bsub
	\item $\TE{A} \TElt \TE{A'}$ if and only if $A \MTle A'$.
	\item For all $N \in \LamI$, we have $\TE{\dirapp N} \TElt \TE N$.
	\esub
\end{lemma}

	\begin{proof}
	Immediate inductions (i) on $A$ and $A'$,
	and (ii) on the head structure of $N$.
	\end{proof}

\newcounter{lem:rred-simul-bred}
\setcounter{lem:rred-simul-bred}{\thetheorem}
\begin{lemma}[Simulation of $\bred$]
\label{lem:rred-simul-bred}
	If $M \bred N$, then $\TE{N} = 
	\bigTElub_{\rt s \TEin \TE M} \supp*{\rs t_{\rt s}}$
	for resource sums $\rs t_{\rt s} \in \frsums[\fv(M)]$
	such that $\forall \rt s \TEin \TE M,\ \rt s \rreds \rs t_{\rt s}$.
\end{lemma}

	\begin{proof}
	One first shows by induction on $P$ that
	for all $P,Q \in \LamI$ and $x \in \fv(P)$,
	$\TE{P \subst Q} = \bigTElub_{\rt s \TEin \TE P}
	\bigTElub_{\rb t \in \TE** Q} \supp*{ \rt s \rsubst t }$.
	Then the proof is an induction on $M \bred N$,
	using the previous equality in the base case.
	See the details of a similar proof in 
	\cite[Lemmas 4.1 and 4.2]{CerdaVaux23}.
	\end{proof}

\newcounter{lem:TE-of-A-is-normal}
\setcounter{lem:TE-of-A-is-normal}{\thetheorem}
\begin{lemma}
\label{lem:TE-of-A-is-normal}
	If $A \in \Appset$ then all $\rt s \TEin \TE A$
	are in $\resource$-nf.
\end{lemma}

	\begin{proof}
	Immediate induction on $A$.
	\end{proof}

Everything is now in place to prove the Commutation \cref{the:commutation}.

\begin{proof}[Proof of \cref{the:commutation}]
	By \cref{rem:fv-of-taylor}, $\fv(\NFT M) = \fv(M) = \fv(\TMT M)$.
	Thus, it is sufficient to prove that, 
	for $\rt t \in \rterms(\fv(M))$:
	$\rt t \TEin \NFT M$ if and only if $\rt t \TEin \TMT M$.
	
	Take $\rt t \TEin \NFT{M}$, \emph{i.e.} $\rt t \in \supp*{\nf(\rt s)}$
	for some $\rt s \TEin \TE M$.
	Then, by \cref{lem:rred-conserv-bred}, there exists
	a reduction $M \breds N$
	such that $\supp*{\nf(\rt s)} \TElt \TE{N}$.
	Hence $\rt t \TEin \TE{N}$ and $\rt t$ is in $\resource$-nf: 
	\cref{lem:normal-approximants} ensures that $\rt t \TEin \TE A$ 
	for some $A\in \Appof{N}=\Appof{M}$ (by \cref{prop:InvariantUnderBeta}).
	This means exactly that $\rt t \TEin \TMT M$.
	
	Conversely, take $\rt t \TEin \TMT M$, \emph{i.e.}
	$\rt t \TEin \TE{A}$ for some $A \in \Appof{M}$.
	By definition there is a reduction $M \breds N$
	such that $A \MTle \dirapp{N}$.
	By \cref{lem:TE-monotone},
	$\rt t \TEin \TE A \TElt \TE{\dirapp N} \TElt \TE N$.
	By iterated applications of \cref{lem:rred-simul-bred}, 
	there is an $\rt s \TEin \TE M$ such that
	$\rt s \rreds \rs t_{\rt s}$ and $\rt t \in \supp*{\rs t_{\rt s}}$.
	By \cref{lem:TE-of-A-is-normal}, $\rt t$ is in $\resource$-nf,
	therefore $\rt t \TEin \supp*{\nf(\rt s)} \TElt \NFT M$.
\end{proof}


\subsection{The \lamI-theory of Ohana trees}

Consider the equivalence $\cM$ on $\LamI$, 
defined by $M =_{\cM} N$ if and only if $\MT{M} = \MT{N}$.
Thanks to \cref{the:commutation},
we are now able to show that this equivalence is a \lamI-theory.

\begin{corollary} \label{cor:MTeqIffNFTeq}
	For $M,N\in\LamI$, $M =_\cM N$ if and only if $\NFT M = \NFT N$.
\end{corollary}
	\begin{proof}
	($\Rightarrow$)
	Immediate by \cref{the:commutation}.
	($\Leftarrow$)
	Take $M,N$ such that $\NFT M = \NFT N$. 
	By \cref{thm:MTbijApp}, it is sufficient to prove 
	$\Appof{M} = \Appof{N}$.
	Take $A\in\Appof{M}$. 
	If $A = \bot_X$, then $A \in \Appof{N}$ is trivial 
		since $\Appof{N}$ is downward closed.
	Otherwise, by hypothesis and \cref{the:commutation},
		$\TE A \TElt \bigTElub_{B \in \Appof N} \TE  B$.
		There is a $B \in \Appof N$ such that
		$\resourceof A \TEin \TE B$,
		where $\resourceof A \TEin \TE A$ is defined by
		\[
		\begin{array}{rcl}
		\resourceof x &\eqdef& x,\\
\resourceof{\lam \vec{x}.yA_1\cdots A_k} &\eqdef&
			\lam \vec{x}.y\resourceof*{A_1}\cdots \resourceof*{A_k},\\	
		\end{array}
		\qquad
		\resourceof*{A} \eqdef
		\begin{cases}
		1_X,&\textrm{if }A = \bot_X,\\
		{[\resourceof A]},&\textrm{otherwise}.
		\end{cases}
		\]
		By induction on $A$, one shows that
		$\resourceof A \TEin \TE B$ implies $A \MTle B$,
		hence $A \in \Appof N$ by downward-closure
		(\cref{lem:DirApp}\eqref{lem:DirApp2}).
	This shows that $\Appof{M} \subseteq \Appof{N}$.
	The converse inclusion is symmetric.
	\end{proof}

\begin{corollary}\label{cor:MisALamITheory} 
	$\cM$ is a \lamI-theory.
\end{corollary}

\begin{proof}
	The relation $=_\cM$ is clearly an equivalence.
	By \cref{prop:InvariantUnderBeta,thm:MTbijApp}, 
	for all $M,N\in\LamI$, $M =_\beta N$ entails $M =_\cM N$.
	
	For compatibility with abstraction, 
	take $M =_\cM N$ and $x \in \fv(M) \cap \fv(N)$. 
	By \cref{cor:MTeqIffNFTeq}, we get 
	$\NFT{\lam x.M} = \lam x.\NFT M = \lam x.\NFT N = \NFT{\lam x.N}$.
	
	For compatibility with application, 
	by \cref{the:rred-norm-confl} observe that for all $M,N \in \LamI$, 
	$\NFT{MN} = \nf(\TE M \TE** N) = \nf(\NFT M \NFT** N)$,
	where $\NFT** N$ is defined by adapting
	\cref{not:support} to sets of resource bags.
	Therefore, $\NFT M = \NFT{M'}$ and $\NFT N = \NFT{N'}$ 
	imply $\NFT{MN} = \NFT{M'N'}$, and
	we can conclude the proof by \cref{cor:MTeqIffNFTeq}.
\end{proof}
